% Vector schematic: preparation above, memory-to-cache-to-core traffic below.
% Icons are drawn locally so the figure stays scalable and matches the thesis.
\begingroup
\resizebox{\textwidth}{!}{%
\begin{tikzpicture}[
 font=\small\sffamily,text=black!80,
 data/.style={draw=figcol,fill=figcol!8,line width=0.65pt},
 mechanism/.style={draw=tabline,fill=neutralcol,line width=0.65pt},
 flow/.style={-{Stealth[length=5pt,width=4pt]},draw=black!65,line width=0.85pt},
 note/.style={font=\footnotesize\sffamily,align=center},
 label/.style={align=center,text width=2.9cm},
 heading/.style={font=\small\sffamily\bfseries,anchor=west}
]
\node[heading] at (-0.9,1.6) {1\quad Prepare addressable weights};

% Serialized tensor file: a folded document with a numerical array and metadata.
\begin{scope}[shift={(0.5,0.25)}]
 \draw[data] (-0.62,-0.72)--(0.62,-0.72)--(0.62,0.42)
   --(0.32,0.72)--(-0.62,0.72)--cycle;
 \draw[draw=figcol,line width=0.65pt] (0.32,0.72)--(0.32,0.42)--(0.62,0.42);
 \foreach \x in {0,1,2} \foreach \y in {0,1}
  \draw[draw=figcol!65,fill=white] (-0.4+\x*0.25,0.02+\y*0.23)
   rectangle ++(0.2,0.18);
 \draw[draw=figcol!65] (-0.4,-0.22)--(0.37,-0.22);
 \draw[draw=figcol!65] (-0.4,-0.42)--(0.12,-0.42);
\end{scope}
\node[label] at (0.5,-0.85) {Model files};
\draw[flow] (1.4,0.25)--(2.9,0.25);

% Loading makes tensor bytes addressable: storage and memory buffer symbols.
\begin{scope}[shift={(4,0.25)}]
 \draw[mechanism,rounded corners=1pt] (-0.9,-0.5) rectangle (-0.3,0.5);
 \foreach \y in {-0.3,0,0.3}
  \draw[draw=black!45] (-0.76,\y)--(-0.44,\y);
 \draw[flow] (-0.15,0)--(0.35,0);
 \draw[mechanism,rounded corners=1pt] (0.5,-0.5) rectangle (1,0.5);
 \foreach \y in {-0.25,0,0.25}
  \draw[draw=black!45] (0.5,\y)--(1,\y);
\end{scope}
\node[label] at (4,-0.85) {Read or map};
\draw[flow] (5.25,0.25)--(6.5,0.25);

% Optional preparation: rearrange a matrix into a kernel-friendly layout.
\begin{scope}[shift={(7.8,0.25)}]
 \foreach \x in {0,1,2} \foreach \y in {0,1,2}
  \draw[mechanism] (-1+\x*0.22,-0.34+\y*0.22) rectangle ++(0.18,0.18);
 \draw[flow] (-0.18,0)--(0.3,0);
 \foreach \x in {0,1,2,3} \foreach \y in {0,1}
  \draw[mechanism] (0.43+\x*0.19,-0.24+\y*0.28) rectangle ++(0.15,0.22);
\end{scope}
\node[label] at (7.8,-0.85) {Convert / pack};
\node[note,text=black!60] at (7.8,-1.35) {optional};
\draw[flow] (9.1,0.25)--(10.5,0.25);

% Runtime representation: overlapping tensor planes, not physical page residency.
\begin{scope}[shift={(11.7,0.25)}]
 \draw[data] (-0.42,-0.28) rectangle (0.88,0.72);
 \draw[data] (-0.65,-0.51) rectangle (0.65,0.49);
 \foreach \x in {1,2,3}
  \draw[draw=figcol!55] (-0.65+\x*0.325,-0.51)--++(0,1);
 \foreach \y in {1,2}
  \draw[draw=figcol!55] (-0.65,-0.51+\y/3)--++(1.3,0);
\end{scope}
\node[label] at (11.7,-0.85) {Runtime tensors};
\draw[draw=tabline,dashed] (-0.9,-1.9)--(13.2,-1.9);

\node[heading] at (-0.9,-2.6) {2\quad Fetch operands during inference};

% Main-memory pages: one contiguous row is selected for the illustrated fill.
\node[label,font=\small\sffamily\bfseries] at (0.5,-3.7) {Backing memory};
\begin{scope}[shift={(0.5,-5.15)}]
 \draw[data,fill=white] (-0.82,-0.78) rectangle (1.28,0.88);
 \draw[data,fill=white] (-1.0,-0.96) rectangle (1.1,0.7);
 \draw[data,fill=white] (-1.18,-1.14) rectangle (0.92,0.52);
 \foreach \x in {0,1,2,3} \foreach \y in {0,1,2,3}
  \draw[draw=figcol!35,fill=figcol!4]
   (-1.03+\x*0.48,-0.98+\y*0.35) rectangle ++(0.43,0.29);
 \foreach \x in {0,1,2,3}
  \draw[data,fill=figcol!30,line width=0.8pt]
   (-1.03+\x*0.48,-0.28) rectangle ++(0.43,0.29);
\end{scope}
\node[note] at (0.5,-6.65) {HBM and/or DDR};

% CPU boundary contains the cache hierarchy and the executing core.
\draw[draw=tabline,rounded corners=5pt,line width=0.8pt]
 (5.3,-3.85) rectangle (13.2,-7.1);
\node[note,anchor=west,font=\footnotesize\sffamily\bfseries]
 at (5.45,-3.6) {CPU (schematic)};

% Cache rows show retained copies; highlight only the illustrated line.
\begin{scope}[shift={(7,-5.15)}]
 \draw[mechanism,rounded corners=3pt] (-1.2,-1.05) rectangle (1.2,0.75);
 \foreach \x in {0,1,2,3} \foreach \y in {0,1,2}
  \draw[draw=tabline,fill=white]
   (-1.01+\x*0.51,-0.82+\y*0.46) rectangle ++(0.45,0.32);
 \foreach \x in {0,1,2,3}
  \draw[data,fill=figcol!30,line width=0.8pt]
   (-1.01+\x*0.51,-0.28) rectangle ++(0.45,0.32);
\end{scope}
\node[note] at (7,-6.65) {CPU caches};

% A cache-line packet on the data path makes the transfer granularity visible.
\draw[flow] (1.8,-5.27)--(5.78,-5.27);
\foreach \x in {0,1,2,3}
 \draw[data,fill=figcol!30,line width=0.8pt]
  (2.9+\x*0.44,-5.43) rectangle ++(0.39,0.32);
\node[note] at (3.75,-4.65) {Cache-line fill};
\node[note,text=black!60] at (3.75,-5.9) {on a miss to memory};

% Core package with pins and a small matrix multiplication kernel.
\begin{scope}[shift={(11.4,-5.15)}]
 \foreach \p in {-0.55,-0.18,0.18,0.55} {
  \draw[draw=black!50,line width=1pt] (\p,-0.9)--(\p,-1.08);
  \draw[draw=black!50,line width=1pt] (\p,0.9)--(\p,1.08);
  \draw[draw=black!50,line width=1pt] (-0.9,\p)--(-1.08,\p);
  \draw[draw=black!50,line width=1pt] (0.9,\p)--(1.08,\p);
 }
 \draw[mechanism,rounded corners=2pt] (-0.9,-0.9) rectangle (0.9,0.9);
 \draw[draw=black!30,fill=white] (-0.72,-0.72) rectangle (0.72,0.72);
 \foreach \x in {0,1} \foreach \y in {0,1} {
  \draw[mechanism] (-0.5+\x*0.19,0.06+\y*0.19) rectangle ++(0.15,0.15);
  \draw[data,fill=figcol!30] (0.16+\x*0.19,0.06+\y*0.19) rectangle ++(0.15,0.15);
 }
 \node[font=\small\sffamily] at (0,0.24) {$\times$};
 \node[note] at (0,-0.34) {Kernel};
\end{scope}
\node[note] at (11.4,-6.65) {CPU core};
\draw[flow] (8.22,-5.27)--(10.3,-5.27);
\node[note] at (9.3,-4.65) {Operands};

% Legend uses swatches, with an additional key for the highlighted data line.
\node[data,rounded corners=1pt,minimum width=0.28cm,
 minimum height=0.22cm,inner sep=0pt] at (-0.55,-7.85) {};
\node[note,anchor=west] at (-0.28,-7.85) {Weight data};
\node[mechanism,rounded corners=1pt,minimum width=0.28cm,
 minimum height=0.22cm,inner sep=0pt] at (3.6,-7.85) {};
\node[note,anchor=west] at (3.87,-7.85) {Preparation / hardware};
\node[data,fill=figcol!30,minimum width=0.28cm,
 minimum height=0.22cm,inner sep=0pt] at (9.3,-7.85) {};
\node[note,anchor=west] at (9.57,-7.85) {Selected weight data};
\end{tikzpicture}%
}
\endgroup
